\section*{Chapter \ref{ch:ml:clustering-regimes-and-anomaly-detection} --- Clustering, Regimes and Anomaly Detection}

\begin{solution}{exo:ml:clustering-regimes-and-anomaly-detection:1}
$\sqrt{2(1 - 0.5)} = 1$. Perfectly correlated names are at distance 0, uncorrelated ones at $\sqrt2 = 1.41$ (and perfectly
anti-correlated at 2).
\end{solution}

\begin{solution}{exo:ml:clustering-regimes-and-anomaly-detection:2}
Durations $1/(1 - p)$: 82 days calm and 24 days stress (the fitted matrix, 0.988 and 0.958 rounded, gives 82.2 and 23.9). The
long-run share of stress is $P_{12}/(P_{12} + P_{21}) = 0.225$, against 0.215 on the test years.
\end{solution}

\begin{solution}{exo:ml:clustering-regimes-and-anomaly-detection:3}
115 legitimate account-days are flagged at 1\%; with all 100 true cases flagged too, precision is $100/215 = 0.47$.
\end{solution}

\begin{solution}{exo:ml:clustering-regimes-and-anomaly-detection:4}
The filter uses only data up to that day, and the day before a switch is a calm day: its return is drawn from the calm
distribution. A switch is announced by the stress days that follow it, which the filter has not seen; its probability on
that day is just the calm-day probability of a switch, about 1\%, pushed around by the day's return.
\end{solution}

\begin{solution}{exo:ml:clustering-regimes-and-anomaly-detection:5}
Every stock loads on the market, so raw correlations are dominated by one common factor: all names look alike, and the
\omterm{def:ml:clustering-regimes-and-anomaly-detection:cluster}{clustering} picks up differences in beta or specific volatility rather than industries. Removing the market leaves the
industry correlations to drive the distances.
\end{solution}

\begin{solution}{exo:ml:clustering-regimes-and-anomaly-detection:6}
The probabilities are smoothed, or the model was fitted on the whole sample, drawdowns included: its states and parameters know
where the drawdowns were, and the smoothed probability rises before each switch by construction (0.35 against 0.02 on calm
days here). Rerun the model as a filter, refitted only on data before each date, and look again.
\end{solution}

\begin{solution}{exo:ml:clustering-regimes-and-anomaly-detection:7}
The test years' log-likelihood is 7\,866.6 under the two-state model and 7\,864.4 under the three-state one: the extra state,
which improved the fit in sample by 5.7, loses 2.2 out of sample. The information criterion's verdict holds.
\end{solution}

\begin{solution}{exo:ml:clustering-regimes-and-anomaly-detection:8}
$\alpha_t(j) = p(x_1,\dots,x_t, s_t = j) = \sum_ip(x_1,\dots,x_{t-1}, s_{t-1} = i)\P(s_t = j\mid s_{t-1} = i)p(x_t\mid s_t = j)$,
using the Markov property of the state and the conditional independence of $x_t$ given $s_t$; that is
$\big(\sum_i\alpha_{t-1}(i)P_{ij}\big)b_j(x_t)$. Dividing by $\sum_j\alpha_t(j) = p(x_1,\dots,x_t)$ gives $\P(s_t = j\mid
x_1,\dots,x_t)$, the filtered probability; the code normalises at every step, which also keeps the numbers from underflowing.
\end{solution}

\begin{solution}{pb:ml:clustering-regimes-and-anomaly-detection:1}
\textbf{Part I.}
\begin{enumerate}
\item With industries at 15\%: adjusted Rand index 0.38 to 0.42 on three months, 0.60 to 0.75 on a year, 0.99 to 1.00 on five
years. At 8\%: 0.03 to 0.06, 0.06, and 0.15 to 0.21.
\item Half-sample agreement over a year: 0.32 to 0.35 (strong industries), 0.06 to 0.08 (weak). On real data there are no
planted labels; stability is the only check.
\item Grouping for risk reports and hierarchical risk parity when stable; not as a trading signal, and not when they do not
survive a split.
\item \omterm{def:ml:clustering-regimes-and-anomaly-detection:cluster}{Hierarchical clustering} on the correlation distance is slightly better here and gives a tree that shows how groups
nest; \omterm{def:ml:clustering-regimes-and-anomaly-detection:cluster}{k-means} needs a fixed number of groups and a Euclidean representation.
\end{enumerate}
\textbf{Part II.}
\begin{enumerate}[resume]
\item Calm volatility 0.80\% and stress 2.60\% (planted 0.8\% and 2.5\%), staying probabilities 0.988 and 0.958 (0.99 and
0.95).
\item The third state adds 5.7 to the in-sample log-likelihood for six parameters, below the information criterion's 23.5, and
loses out of sample (exercise~7).
\item Filtered 95.9\%, smoothed 97.7\%, Viterbi 97.5\%, volatility threshold 89.8\%.
\item Median one day for the filter (mean 1.5), zero for the smoother (mean 0.8); the volatility threshold's mean is 3.8 and it
misses five of 20 spells.
\end{enumerate}
\textbf{Part III.}
\begin{enumerate}[resume]
\item Filtered 0.032 (calm-day average 0.041); smoothed 0.353 (calm-day average 0.019).
\item Sharpe ratio 1.21 with the smoothed probability against 1.00 with the filtered one (buy and hold 0.91).
\item From filtered probabilities, with parameters fitted on data before each date.
\item In any statistic computed with a full-sample fit: normalisations, regime labels, detrending, outlier cleaning, and
Kalman smoothers used where a filter was meant (Book~4, chapter~19).
\end{enumerate}
\textbf{Part IV.}
\begin{enumerate}[resume]
\item \emph{Which regime are we in?} The filtered regime is right on 95.9\% of days and sees a switch into stress after a median
of one day; the smoothed one is right on 97.7\% and sees it on the day, because it knows the future. At a 1\% false-positive
rate the \omterm{def:ml:clustering-regimes-and-anomaly-detection:anomaly}{isolation forest}'s precision is 0.47 (recall 1.00), the best rule's 0.40 (0.77) and the \omterm{def:ml:cross-sectional-deep-models:ae}{autoencoder}'s 0.22 (0.32).
\item Describing the present state for risk scaling and for conditioning other models, quickly and honestly, not predicting
switches.
\item No: as a second net, reviewed, with the rules as the primary detectors whose logic can be explained to a regulator.
\item By the features that made each alarm unusual, grouped by account type, against the rules' verdicts, with outcomes fed
back as labels.
\item Its filtered probabilities' calibration (the frequency of stress when it says 0.8), parameter drift across refits, and
out-of-sample likelihood.
\item The tree sets the allocation's hierarchy; noise clusters make the weights jump with each refit and add turnover without
diversification.
\item Confirmed cases from investigations: slow, few, biased towards what the rules already find.
\item What yesterday's data say, and no more.
\end{enumerate}
\end{solution}

\begin{solution}{iq:ml:clustering-regimes-and-anomaly-detection:1}
Filtered: the state's probability given observations up to $t$, available at $t$. Smoothed: given all observations, including
after $t$; better for describing history, look-ahead if used for decisions.
\iqlookfor{the information sets and the look-ahead implication.}
\end{solution}

\begin{solution}{iq:ml:clustering-regimes-and-anomaly-detection:2}
Remove the market (and possibly styles), compute correlations over a window long enough for the effects, cluster on the
correlation distance, and check stability across sub-samples and against known classifications.
\iqlookfor{factor removal, a distance, and a stability check.}
\end{solution}

\begin{solution}{iq:ml:clustering-regimes-and-anomaly-detection:3}
EM for HMMs: the E-step runs forward-backward to get each state's and each transition's posterior probability; the M-step
re-estimates transition probabilities and emission parameters as posterior-weighted averages; repeat until the likelihood stops
rising.
\iqlookfor{forward-backward, posterior weights, re-estimation, likelihood monotonicity.}
\end{solution}

\begin{solution}{iq:ml:clustering-regimes-and-anomaly-detection:4}
Random trees split on random features at random values; anomalies need fewer splits to be isolated. Blind spots: anomalies that
are not rare (clusters of manipulators), anomalies visible only in combinations the random splits rarely separate, and
irrelevant features that add noise.
\iqlookfor{the mechanism and at least two blind spots.}
\end{solution}

\begin{solution}{iq:ml:clustering-regimes-and-anomaly-detection:5}
Whether the regime probabilities were filtered and fitted only on past data; whether the number of regimes and parameters were
chosen on the whole sample; the result's significance; and its stability across sub-periods.
\iqlookfor{smoothing and full-sample fitting as the first suspects.}
\end{solution}

\begin{solution}{iq:ml:clustering-regimes-and-anomaly-detection:6}
Rules for known patterns plus unsupervised scores on engineered features; thresholds set by reviewer capacity; evaluation on
planted scenarios and look-alikes, on the few confirmed cases, and by reviewers' dispositions over time; an audit trail of every
alarm and decision.
\iqlookfor{rules and scores together, capacity-based thresholds, planted-case evaluation, feedback loop.}
\end{solution}
